% ==================== 三、问题一 ====================
\section{问题一：典型日确定性储能调度}

\subsection{问题分析}

问题一给定某典型日的电价、小区负载与光伏发电预测功率，要求制定使购电费最小的调度策略。
由于该日全部数据已知，这是一个\textbf{确定性优化问题}。

其经济本质是两条：
\begin{enumerate}[itemsep=2pt,topsep=3pt]
  \item \textbf{光伏余电消纳}：光伏出力超过负载时，若不储能则被迫弃光；
        储能可把余电存起来供后续使用，减少购电。
  \item \textbf{峰谷价差套利}：在低价时段多购电存入电池，在高价时段放电使用，
        赚取价差。能否套利取决于价差是否超过往返效率损失。
\end{enumerate}

\subsection{符号说明}

本问使用的全部符号已在 \S\ref{sec:notation} 表 \ref{tab:notation} 统一说明，此处不再重复。
下面直接给出模型。

\subsection{模型建立}
\label{sec:q1-model}

以全天购电费用最小为目标：

\begin{equation}
\min\ J_1=\sum_{t=1}^{144} p_t g_t
\label{eq:q1-obj}
\end{equation}

约束条件包括五组：

\textbf{(1) 功率平衡}（供给侧 $=$ 需求侧）：
\begin{equation}
g_t+V_t+d_t=L_t+c_t+w_t
\label{eq:q1-balance}
\end{equation}
其中 $w_t$ 为弃光量，故外购电与光伏之和，加上电池放电，恰好满足负载与电池充电。

\textbf{(2) 弃光非负}：
\begin{equation}
w_t\ge 0
\end{equation}
弃光量只设非负约束，\textbf{不显式添加} $w_t\le V_t$ 的上界。由等式平衡
$w_t=g_t+d_t-c_t-(L_t-V_t)$；目标只惩罚购电量，若最优解出现 $w_t>V_t$
（即净供给超过负载），令 $g_t\leftarrow g_t-\delta$、$w_t\leftarrow w_t-\delta$
可保持平衡不变而费用严格下降（对任意 $p_t>0$），与最优性矛盾。
因此最优解自动满足 $w_t\le V_t$，该上界冗余；省略它与第二、三问对松弛变量的
口径保持一致。本节末的约束核验同时报告“弃光非负”（属模型约束）与
“弃光未超光伏出力”（属本实例的数值观察，本实例弃光量为零）。

\textbf{(3) 储能电量递推}（注意充电乘效率、放电除效率）：
\begin{equation}
E_t=E_{t-1}+0.9\,c_t-\frac{d_t}{0.9},\qquad t=1,\dots,144
\label{eq:q1-soc}
\end{equation}

\textbf{(4) 储能容量与功率约束}：
\begin{equation}
1200\le E_t\le 10800,\qquad 0\le c_t\le \frac{5000}{6},\qquad 0\le d_t\le \frac{5000}{6}
\end{equation}

\textbf{(5) 非负与循环条件}：
\begin{equation}
g_t\ge 0,\qquad E_0=E_{144}=6000
\label{eq:q1-cycle}
\end{equation}

模型共含 $5\times 144=720$ 个决策变量，全部约束为线性，
构成一个稀疏线性规划（LP），调用 \texttt{scipy.optimize.linprog}（HiGHS 内点/单纯形）求解。

\subsection{关键性质：LP 松弛的紧性证明}
\label{sec:q1-tight}

物理上储能不能同时充电与放电，即要求 $c_t\cdot d_t=0$。
但引入这类互斥条件需要二元变量，问题变为混合整数规划（MILP），规模与计算量显著上升。
我们证明：\textbf{在本模型结构下，最优解自动满足互斥，无需显式约束}。

\begin{quote}
\textbf{命题。} 设某时段同时有 $c_t>0$ 与 $d_t>0$。取
\[
\varepsilon=\min\Bigl\{c_t,\ 0.81\,d_t,\ \frac{g_t}{0.19}\Bigr\}>0,
\]
并作替换 $c_t\leftarrow c_t-\varepsilon$、$d_t\leftarrow d_t-0.81\varepsilon$，其余变量不变。则
\begin{enumerate}[label=(\roman*),itemsep=1pt,topsep=2pt]
  \item \textbf{储电量轨迹与全部 SOC 约束完全不变}，因为
    \[
    0.9(c_t-\varepsilon)-\frac{d_t-0.81\varepsilon}{0.9}
    =\Bigl(0.9c_t-\frac{d_t}{0.9}\Bigr)-0.9\varepsilon+0.9\varepsilon
    =\Bigl(0.9c_t-\frac{d_t}{0.9}\Bigr);
    \]
  \item \textbf{充电、放电功率均不增加}（故功率约束仍满足）；
  \item 由功率平衡 \eqref{eq:q1-balance}，$g_t$ 变为 $g_t-0.19\varepsilon<g_t$，
        即\textbf{购电量严格下降}，费用随之严格下降。
\end{enumerate}
这与最优性矛盾。故最优解中必不存在“该时段有购电且同时充放电”的情形。
\end{quote}

\textbf{经济解释}：往返效率 $0.81<1$ 意味着“买 1 度电只有 0.81 度能存进电池再取出”。
当该时段仍在从外网购电时，把充放电的这部分对冲掉，等价于\textbf{少买一点电、同时少损耗一点}，
而负载供应不受影响——因此必然更省钱。

\textbf{验证}：本实例中 LP 最优解的 144 个时段\textbf{全部满足 $c_td_t=0$}，
无需任何后处理，故 LP 松弛是紧的。求解 MILP 作为对照，两者目标值完全相同
（$35\,126.95$ 元，差 $0.000000$），从数值上证实了上述证明。

\subsection{求解结果}

\begin{table}[H]
\centering
\caption{问题一求解结果对比}
\label{tab:q1-result}
\begin{tabular}{@{}l r r@{}}
\toprule
方案 & {购电费（元）} & {相对基线} \\
\midrule
无储能基线 & 48052.05 & — \\
\textbf{储能在位、优化调度} & \bfseries 35126.95 & \bfseries $-26.90\%$ \\
\bottomrule
\end{tabular}
\end{table}

经济性校验：充电加权均价 $0.4930$ 元/kWh，放电加权均价 $1.1376$ 元/kWh，
比值 $2.31$，高于盈亏平衡点 $1/0.81=1.235$，故储能套利有正收益。

\subsection{收益来源分解}

把优化结果与无储能基线逐项对比，可把 $12\,925.10$ 元的节省\textbf{严格分解}为两部分。
由逐时段功率平衡 $g_t=L_t+c_t-d_t-V_t+w_t$，
基线取 $c=d=0$ 且\ $w_t=\max(V_t-L_t,0)$，可得\textbf{严格恒等式}

\begin{equation}
\sum_t p_t\max(L_t-V_t,0)-\sum_t p_tg_t
=\sum_t p_t\max(V_t-L_t,0)+\sum_t p_td_t-\sum_t p_tc_t.
\label{eq:q1-decomp}
\end{equation}

即“总节省 $=$ 弃光消除 $+$ 净套利”，左端为基线购电费与优化购电费之差，
右端两项分别对应\textbf{被吸收的光伏余电价值}与\textbf{储能低买高卖的净收益}。
两者的经济含义见下表，\textbf{各项之和恒等于总节省，无重复、无遗漏}。

\begin{table}[H]
\centering
\small
\caption{问题一收益来源分解（严格恒等式 \eqref{eq:q1-decomp}）}
\label{tab:q1-gain}
\begin{tabularx}{\textwidth}{@{}l r r X@{}}
\toprule
收益来源 & {金额（元）} & {占基线} & 说明 \\
\midrule
光伏余电消纳 & 4037.75 & 8.40\% & 基线中被弃掉的 $6247.96\kw$ 全部被电池吸收 \\
峰谷价差套利 & 8887.35 & 18.50\% & 低价购电存入、高价放出，已扣除充电成本 \\
\midrule
\textbf{合计} & \bfseries 12925.10 & \bfseries 26.90\% & \\
\bottomrule
\end{tabularx}
\end{table}

可见\textbf{套利是主要来源（约 69\%），消纳次之（约 31\%）}。
这是因为本例中光伏余电仅 $6247.96\kw$（占光伏总量 $11.3\%$），
而峰谷价差最高达 $1.40/0.37\approx 3.8$ 倍，远超往返效率盈亏平衡点 $1/0.81=1.235$，
故储能的边际价值主要体现在跨时段搬移电量而非吸收余电。

\subsection{电量平衡与调度解释}

\begin{table}[H]
\centering
\caption{问题一全天电量平衡}
\begin{tabular}{@{}l r l r@{}}
\toprule
项目 & {电量（kWh）} & 项目 & {电量（kWh）} \\
\midrule
外网购电   & 59482.70 & 小区负载     & 111024.81 \\
光伏发电   & 55482.84 & 电池充电     & 20740.67 \\
电池放电   & 16799.94 & 弃光         & \bfseries 0.00 \\
         &          & 电池往返损失 & 3940.73 \\
\bottomrule
\end{tabular}
\end{table}

\textbf{为何弃光为 0？}光伏总量 $55\,482.84\kw$ 全部被利用，
其中一部分直接供给负载、一部分存入电池。这说明在本日的“光伏—负载—电价”组合下，
电池容量足以吸收全部余电，储能的边际价值主要体现在套利而非消纳。

\textbf{为何充放电发生在这些时段？}附件 1 的电价并非简单的“两段式”，
而是呈现\textbf{三段特征}：

\begin{table}[H]
\centering
\caption{典型日电价的时段结构}
\begin{tabular}{@{}l l l@{}}
\toprule
时段 & 电价水平（元/kWh） & 调度行为 \\
\midrule
凌晨 $0$--$6$ 时   & 低谷 $0.37$--$0.44$ & 大量充电（$0$--$4$ 时充 $4500\kw$）\\
上午 $6$--$9$ 时    & \textbf{尖峰}，最高 $1.01$  & 放电供负载 \\
中午 $10$--$16$ 时   & 回落 $0.41$--$0.53$   & 光伏余电 + 补充电 \\
傍晚 $18$--$20$ 时   & \textbf{全天最高} $1.35$--$1.40$ & 放电（$16$--$20$ 时放 $5780\kw$）\\
\bottomrule
\end{tabular}
\end{table}

调度的全部逻辑围绕这三段展开：在凌晨与中午的低价窗口充电，
在上午与傍晚的两个高价窗口放电，从而最大化价差收益。

\subsection{对偶分析}

对功率平衡约束 \eqref{eq:q1-balance} 取对偶变量（影子价格）$\lambda_t$，
其经济含义为“第 $t$ 段若能多供 1 kWh，购电费能下降多少”。
由最优性条件，$\lambda_t\le p_t$ 时该段不购电（$g_t=0$）；
$\lambda_t=p_t$ 时该段购电为正。对偶量可用来识别“哪些时段的容量瓶颈最值钱”，
具体数值已保存于结果文件的 \texttt{duals.npz} 中。

\begin{figure}[htbp]
\centering
\includegraphics[width=0.95\textwidth]{01_q1_dispatch.pdf}
\caption{典型日确定性调度结果。四个面板共享时间轴，依次展示电价、供需与购电功率、
储能充放电功率及储电量；水平虚线为 1200 与 10800 kWh 的储电量边界，
浅色背景仅辅助定位早晚时段。}
\label{fig:q1-dispatch}
\end{figure}

储电量始终位于规定区间内，且日初、日末均为 6000 kWh，
因此该收益并非通过消耗日初库存获得。

\begin{figure}[htbp]
\centering
\includegraphics[width=0.95\textwidth]{02_q1_savings.pdf}
\caption{储能调度相对无储能基线的购电差异与费用桥接。左图为优化购电功率
减去无储能购电功率；右图按同一时段电价分别累计少购所减少的费用和增购所增加的费用。}
\label{fig:q1-savings}
\end{figure}

图 \ref{fig:q1-savings} 给出了节省的账面来源：储能使部分时段少购电，
对应费用减少 $19\,111.64$ 元；同时在另一些时段增购电，对应费用增加 $6\,186.54$ 元。
两者抵消后净节省 $12\,925.10$ 元，降幅为 $26.90\%$。
这一分解是逐时段购电费差额的恒等核算，图 \ref{fig:q1-gain} 的
$4037.75$ 元与 $8887.35$ 元则是按能量来源的另一组恒等分解，
两种口径各自自洽，不互相替代，也不把同一节省重复归因为独立效应。

\begin{figure}[H]
\centering
\includegraphics[width=0.95\textwidth]{q1_comparison.pdf}
\caption{问题一收益来源分解：弃光消除项与净套利项之和恒等于总节省
（式 \eqref{eq:q1-decomp}）}
\label{fig:q1-gain}
\end{figure}

\subsection{题目指定表格}

\begin{table}[H]
\centering
\caption{表 1\quad 微网在指定时间段的购电量及全天购电量与购电费}
\begin{tabular}{@{}l r r@{}}
\toprule
时间段 & {购电量（kWh）} & {该时段电价（元/kWh）} \\
\midrule
10:00--10:10 & 0.00   & 0.9919 \\
12:00--12:10 & 480.41 & 0.4411 \\
14:00--14:10 & 0.00   & 1.0002 \\
16:00--16:10 & 445.43 & 0.8271 \\
18:00--18:10 & 531.89 & 1.2561 \\
20:00--20:10 & 0.00   & 1.3513 \\
\midrule
\textbf{全天购电量} & \bfseries 59482.70 & \\
\textbf{全天购电费} & \bfseries 35126.95 元 & \\
\bottomrule
\end{tabular}
\end{table}

\begin{table}[H]
\centering
\caption{表 2\quad 储能设备在指定时间段的充放电量及 $0$:00 与 $24$:00 的储电量}
\begin{tabular}{@{}l r r@{}}
\toprule
时间段 & {充电量（kWh）} & {放电量（kWh）} \\
\midrule
0:00--4:00   & 4500.00 & 0.00 \\
4:00--8:00   & 833.33  & 6365.84 \\
8:00--12:00  & 4787.96 & 1703.00 \\
12:00--16:00 & 5286.04 & 91.10 \\
16:00--20:00 & 0.00    & 5780.13 \\
20:00--24:00 & 5333.33 & 2859.87 \\
\midrule
$0$:00 储电量  & \multicolumn{2}{c}{6000.00 kWh} \\
$24$:00 储电量 & \multicolumn{2}{c}{6000.00 kWh} \\
\bottomrule
\end{tabular}
\end{table}
